\begin{afterword}
\summaryhead

An \textsc{X-Men} trailer says that every human carries the genetic code for mutation, and the mutant Storm throws lightning. The author argues that an animal gaining the organs and control this needs in one generation, by mutation, would falsify evolutionary theory; a theory that covered Storm would explain anything. The comics use ``evolution'' and ``mutation'' only to place themselves in the literary genre of science.

Many people who believe in evolution, the post says, do the same. They do not know what the theory permits and forbids, speak of the next step in human evolution as if selection had a plan, and call chip designs or corporate splits ``evolution.'' Probably most of them say ``because of evolution'' to belong to the scientific in-crowd. Critics of strongly superhuman AI likewise dismiss it as pseudoscience by comparing it to apocalyptic literature, not by any model. The post ends by asking readers which experiences their proudest scientific beliefs forbid; anything else is ``probably passwords or attire.''

\responsehead

The post's opening argument is sound. Storm's powers would need several complex adaptations at once, and evolutionary theory forbids their arising from one mutation in one generation. A theory extended to cover her would forbid nothing. The rest of the post applies the same test to people, and there it fails.

The main charge is about motives, and it has no evidence. ``Probably an actual majority'' of believers in evolution, the post says, use ``because of evolution'' to belong to a tribe, and would say ``because of intelligent design'' as cheerfully if the in-crowd did. The post's own test, asking what a belief forbids, can show that a person's belief forbids nothing. It cannot show why the person holds it. The post supplies the motive itself. This is the gap already noted at ``Belief as Attire'': no way to tell a belief worn for belonging from one held on grounds the group also holds.

One of the examples misreads a word. People who call a better chip design ``evolution'' are using the word's ordinary sense, which Merriam-Webster gives as a process of change toward a better or more complex state, with the example sentence ``an important step in the evolution of computers.'' They are not claiming that evolutionary biology covers chip design, so their usage does not show that they think the theory covers everything.

The paragraph on artificial intelligence applies the test to opponents who are not named. No critic is quoted. The post faults them for having no model, and offers none for the outcome they doubt. The post it links to, on nanotechnology, at least points to Drexler's reasoning ``from known science,'' and still calls the conclusion ``merely a best guess.'' Here the reader is told that science ``is about models'' and is shown only one side's lack of a model.

In short: a sound argument about a comic-book mutant, extended without evidence to the motives of most believers in evolution, and to unnamed critics of AI who are asked for a model that the post itself does not give.
\end{afterword}
